\documentclass[aps,preprint,superscriptaddress,longbibliography]{revtex4-2}

\usepackage{amsmath,amssymb}
\usepackage{hyperref}
\usepackage{booktabs}
\usepackage{siunitx}

\begin{document}

\title{CRSM and Tetrahedral Dielectric Lattice Resonances: A Hypothesis and Reproducibility Protocol}

\author{Devin Phillip Davis}
\affiliation{Agile Defense Systems LLC, Kentucky, USA}
\affiliation{OSIRIS Quantum Discovery Program}

\date{Draft: \today}

\begin{abstract}
This document defines a testable research program for the proposed Coaxial
Resonance State Manifold (CRSM) and related tetrahedral dielectric-lattice
hypotheses. It is a hypothesis and reproducibility draft, not a report of
confirmed anomalous vacuum energy, propulsion, quantum entropy suppression,
or a replacement for dark matter. Numerical claims in earlier CRSM manuscript
versions, including an entropy-suppression score and a fixed physical lock
angle, are not treated here as verified results. The existing OSIRIS claims
register records \texttt{THETA\_LOCK} as \texttt{REFUTED} and
\texttt{LAMBDA\_PHI} as \texttt{NOT\_MEASURED}; this draft does not alter
those records. We specify the mathematical, statistical, and experimental
artifacts required before any claim can be evaluated for support.
\end{abstract}

\maketitle

\section{Status, scope, and provenance}

\textbf{Status: hypothesis/protocol draft.} No new experiment, hardware job,
numerical reproduction, independent review, or pre-registration is claimed
to have occurred by virtue of this document.

The CRSM label refers to a proposed model family involving dielectric
structures, geometric constraints, and field configurations. The existence
of a useful mathematical parameterization would not by itself establish a
new vacuum mechanism or physical effect. Each claim must be stated with an
observable, units, a null model, an analysis method, and a reproducible
source artifact.

The related historical OSIRIS manuscript
\texttt{papers/crsm\_nature\_physics\_2026.tex} is preserved unchanged.
This draft is a separate artifact and should not be interpreted as a
replacement or as byte-identical to prior manuscript text. The accompanying
repository document \texttt{docs/research/crsm\_claim\_audit.md} maps
important claims to current evidence status and missing evidence.

\section{Claims and current evidence status}

\begin{table}[h]
\centering
\caption{Initial evidence posture for major CRSM claims. These are not new
claim-register transitions.}
\begin{tabular}{p{0.39\linewidth}p{0.18\linewidth}p{0.34\linewidth}}
\toprule
Claim & Current treatment & Evidence needed \\
\midrule
A physical resonance at $51.843^\circ$ & Existing status: \texttt{REFUTED} & Independent, pre-registered measurement with calibrated controls \\
A universal $\Lambda_\Phi$ coupling & Existing status: \texttt{NOT\_MEASURED} & Dimensionally consistent derivation and traceable measurement \\
Quantum entropy suppression & Unverified here & Raw job data, dependence analysis, fixed statistics, replication \\
Casimir or vacuum-energy anomaly & Hypothesis & Correct geometry, boundary conditions, uncertainty budget, conventional baseline \\
Propulsion or mass reduction & Hypothesis & Energy/momentum balance and direct calibrated force measurement \\
Galaxy rotation-curve alternative & Hypothesis & Versioned data and fair, held-out model comparison \\
\bottomrule
\end{tabular}
\end{table}

\section{Geometry must be specified before interpretation}

A tetrahedral construction requires explicit three-dimensional vertices and
a statement of which geometric angle is intended. Common tetrahedral angles
are not interchangeable. In particular,
\begin{equation}
\arccos(1/3) \approx 70.5288^\circ,
\end{equation}
which is not $51.843^\circ$. The latter is approximately
$\arctan(14/11)$ and must not be described as the result of the former
identity. Any proposed lock or resonance angle must be derived from an
explicit geometry and then measured independently; selecting a value from
a geometric coincidence does not establish a physical resonance.

For a proposed lattice, the authors should provide coordinates
$\{\mathbf{v}_i\}$, edge lengths, boundary conditions, orientation
conventions, and a reproducible calculation of any predicted eigenmodes.
A schematic or a fitted angle alone is insufficient.

\section{Dimensional and field-theory requirements}

The ideal parallel-plate Casimir energy per unit area is commonly written
\begin{equation}
\frac{E_{\mathrm{Cas}}}{A}
= -\frac{\pi^2\hbar c}{720d^3},
\end{equation}
while the corresponding force per unit area is
\begin{equation}
\frac{F_{\mathrm{Cas}}}{A}
= -\frac{\pi^2\hbar c}{240d^4}.
\end{equation}
These expressions refer to idealized parallel conducting plates at
separation $d$. A proposed dielectric cavity requires its own geometry,
material response, mode spectrum, and boundary conditions. A quantity
proportional to $d^{-4}$ must not be labeled energy density without a
definition and dimensional derivation appropriate to the model.

A parameter with units of inverse seconds cannot be inserted directly into
a dimensionless multiplicative correction unless a reference timescale or
other dimensionally consistent definition is provided. The numerical value
$2.176435\times10^{-8}$ is close to the Planck mass expressed in kilograms;
that numerical match does not establish a coupling rate in
$\mathrm{s}^{-1}$.

For electromagnetic energy transfer, use the oriented Poynting flux
\begin{equation}
P_{\partial V}(t)=\oint_{\partial V}
\mathbf{S}(\mathbf{x},t)\cdot\hat{\mathbf{n}}\,dA,
\qquad
\mathbf{S}=\mathbf{E}\times\mathbf{H},
\end{equation}
with a defined control volume, surface orientation, time averaging, and
energy accounting. An integral of $|\mathbf{E}\times\mathbf{B}|$ is not
equivalent to a signed net power flux. A flux calculation alone does not
demonstrate propulsion, mass reduction, energy creation, or a spacetime
metric perturbation. Any relativistic interpretation needs an explicit
derivation from a stated stress-energy tensor and field equations.

\section{Quantum-device entropy claim: minimum evidence}

Any previously reported entropy-suppression statistic, including a large
$Z$-score or small $p$-value, remains unverified in this draft. The
minimum reproducibility package includes:
\begin{enumerate}
  \item immutable per-job raw results or counts, job identifiers, backend,
  timestamps, circuit hashes, shot counts, and calibration metadata;
  \item the exact data-selection history, exclusion log, analysis commit,
  dependency lock, random seeds, and executable command;
  \item a pre-specified sampling unit and treatment of dependence between
  jobs, circuits, calibration windows, and repeated runs;
  \item a fixed entropy estimator, baseline definition, effect size,
  uncertainty interval, primary statistic, and multiplicity correction;
  \item all controls, all conditions, and all null or negative results;
  \item an independent reproduction from the immutable artifacts and a
  separately collected replication under a frozen analysis.
\end{enumerate}

A job count is not automatically an independent sample size. A small
$p$-value, even when correctly computed, is not by itself evidence for a
particular physical mechanism. Exploratory analyses must be labeled as
exploratory and may not be retroactively described as confirmatory.

\section{Proposed experimental program}

\subsection{Microwave or dielectric cavity}

Specify the cavity geometry, dielectric properties, mode spectrum,
instrument calibration, thermal conditions, and expected conventional
response before collecting outcome data. Compare the proposed model with
conventional electromagnetic calculations and include sham and negative
controls. Publish uncertainties and all deviations from the frozen
protocol.

\subsection{Optical interferometry}

Define the observable, phase-stability requirements, noise model, blind or
randomized conditions where practical, and a pre-specified decision rule.
Distinguish instrument drift and environmental coupling from a model-specific
prediction. No anomalous result should be claimed without independent
calibration and replication.

\subsection{Galaxy rotation curves}

Use a named, versioned public dataset and preserve source hashes. Freeze
sample selection, distance and inclination treatment, baryonic components,
uncertainties, likelihood, priors, and held-out evaluation before fitting.
Compare CRSM with established alternatives using the same observations and
uncertainty assumptions. Report model complexity, residual diagnostics,
uncertainty intervals, sensitivity analyses, and held-out predictive
performance. A visual match to selected curves is not sufficient evidence
of superiority.

\section{Analysis and reporting policy}

Before confirmatory data analysis, freeze the hypothesis, sampling unit,
primary endpoint, null model, test statistic, exclusions, stopping rule,
multiplicity policy, and failure criteria. Preserve raw data immutably and
record hashes for all derived artifacts. Report effect sizes and uncertainty
intervals, not only significance levels. Publish null outcomes and all
protocol deviations.

Software tests can establish specified properties of code; they do not by
themselves validate a physical model. Simulations are conditional on their
assumptions and do not count as physical measurements. Configured constants
are not measured quantities. Every claim should be linked to a stable source
artifact and a reviewed evidence status.

\section{Limitations and decision rule}

This document supplies a protocol framework, not results. The equations and
tests above identify requirements for a later investigation; they do not
establish that CRSM is physically correct or that an anomaly exists. Claims
remain hypotheses until data provenance, dimensional consistency,
appropriate controls, independent analysis, and replication support the
specific statement being made. Existing negative and artifact findings must
remain visible in the research lineage.

\section*{Data and code availability}

No new dataset, experiment, or analysis output is released with this draft.
A future result must identify its immutable data artifacts, source commit,
environment, exact reproduction commands, output hashes, and review record.
The manuscript should not be assigned a publication DOI or described as
peer-reviewed until those events actually occur.

\end{document}
